\documentclass[11pt,a4paper]{article}
\usepackage[utf8]{inputenc}
\usepackage[T1]{fontenc}
\usepackage{lmodern}
\usepackage[margin=2.4cm]{geometry}
\usepackage{amsmath,amssymb}
\usepackage{booktabs}
\usepackage{xcolor}
\usepackage{hyperref}
\hypersetup{colorlinks=true,linkcolor=blue!50!black,urlcolor=blue!50!black}

\title{Nonlinear dynamics: Newmark scheme and Newton--Raphson}
\author{Dorian Nezzar}
\date{December 2023, revised \today}

\begin{document}
\maketitle

\section{Equation of motion}
For one degree of freedom, with internal forces depending on the solution, the semi-discrete
equation of motion at time step $n+1$ reads
\begin{equation}
  M\ddot u^{n+1} + C\dot u^{n+1} + f^{\text{int}}(u^{n+1}, \boldsymbol\alpha^{n}) - f^{\text{ext}}(u^{n+1},t^{n+1})
  = r(u^{n+1}) = 0 ,
  \label{eq:motion}
\end{equation}
where $\boldsymbol\alpha^n$ denotes the internal (history) variables, known at the end of step $n$
and frozen during the solution of step $n+1$ (for a damage model: the damage $d$ reached so far).

\section{Newmark $\beta$ scheme in predictor--corrector form}
For a time step $\Delta t = t^{n+1}-t^n$ (Belytschko), the displacement and velocity at $n+1$ are
\begin{align}
  u^{n+1} &= \tilde u^{n+1} + \beta\Delta t^2\,\ddot u^{n+1}, &
  \tilde u^{n+1} &= u^n + \Delta t\,\dot u^n + \frac{\Delta t^2}{2}(1-2\beta)\,\ddot u^n, \label{eq:pred_u}\\
  \dot u^{n+1} &= \tilde{\dot u}^{n+1} + \gamma\Delta t\,\ddot u^{n+1}, &
  \tilde{\dot u}^{n+1} &= \dot u^n + (1-\gamma)\Delta t\,\ddot u^n . \label{eq:pred_v}
\end{align}
Inverting \eqref{eq:pred_u} and substituting in \eqref{eq:pred_v}:
\begin{equation}
  \ddot u^{n+1} = \frac{1}{\beta\Delta t^2}\bigl(u^{n+1}-\tilde u^{n+1}\bigr),\qquad
  \dot u^{n+1} = \tilde{\dot u}^{n+1} + \frac{\gamma}{\beta\Delta t}\bigl(u^{n+1}-\tilde u^{n+1}\bigr).
  \label{eq:acc_vel}
\end{equation}
Both \eqref{eq:acc_vel} are substituted into \eqref{eq:motion}, which becomes a nonlinear
equation in the single unknown $u^{n+1}$:
\begin{equation}
  r(u^{n+1}) = \Bigl(\frac{M}{\beta\Delta t^2} + \frac{\gamma\,C}{\beta\Delta t}\Bigr)\bigl(u^{n+1}-\tilde u^{n+1}\bigr)
  + C\,\tilde{\dot u}^{n+1} + f^{\text{int}}(u^{n+1}) - f^{\text{ext}}(u^{n+1},t^{n+1}) = 0 .
  \label{eq:residual}
\end{equation}
With $C = 0$ the damping terms vanish and only the inertia term $M(u^{n+1}-\tilde u^{n+1})/(\beta\Delta t^2)$
is added to the static residual.

\paragraph{Properties.}
$\beta = \tfrac14$, $\gamma = \tfrac12$: average acceleration, second-order accurate,
unconditionally stable, no numerical damping. For $\gamma > \tfrac12$ the scheme is first-order
accurate and dissipative (algorithmic damping growing with $\gamma-\tfrac12$ and with
$\omega\Delta t$); unconditional stability then requires $\beta \ge (\gamma+\tfrac12)^2/4$.

\section{Newton--Raphson}
The residual \eqref{eq:residual} is linearised around the iterate $u_i^{n+1}$ by a Taylor
expansion:
\begin{equation}
  r(u_{i+1}^{n+1}) \approx r(u_i^{n+1}) + \frac{\partial r}{\partial u}\Big|_{u_i^{n+1}}\Delta u_i^{n+1} = 0,
  \qquad \Delta u_i^{n+1} = u_{i+1}^{n+1} - u_i^{n+1},
  \label{eq:taylor}
\end{equation}
with the consistent tangent obtained by differentiating \eqref{eq:residual}:
\begin{equation}
  \frac{\partial r}{\partial u} = \frac{M}{\beta\Delta t^2} + \frac{\gamma\,C}{\beta\Delta t}
  + \frac{\partial f^{\text{int}}}{\partial u} - \frac{\partial f^{\text{ext}}}{\partial u} .
  \label{eq:tangent}
\end{equation}
The last term vanishes when the external force does not depend on the displacement (dead
load, imposed displacement), which is the usual case. Solving \eqref{eq:taylor} for the
increment and updating:
\begin{equation}
  \Delta u_i^{n+1} = -\Bigl(\frac{\partial r}{\partial u}\Big|_{u_i^{n+1}}\Bigr)^{-1} r(u_i^{n+1}),
  \qquad
  u_{i+1}^{n+1} = u_i^{n+1} + \Delta u_i^{n+1}.
  \label{eq:update}
\end{equation}
For several degrees of freedom, $\partial r/\partial u$ is the effective tangent matrix
$\mathbf K_{\text{eff}} = \mathbf M/(\beta\Delta t^2) + \gamma\mathbf C/(\beta\Delta t) + \mathbf K_t$
and \eqref{eq:update} is a linear solve.

\section{Algorithm for one time step}
\begin{enumerate}
  \item Predictors $\tilde u^{n+1}$, $\tilde{\dot u}^{n+1}$ from \eqref{eq:pred_u}--\eqref{eq:pred_v}.
  \item Initial iterate $u_0^{n+1} = \tilde u^{n+1}$; internal variables frozen at $\boldsymbol\alpha^n$.
  \item Iterate \eqref{eq:update} until $|r(u_i^{n+1})| \le \varepsilon$ (absolute tolerance) or
        $\|\Delta u_i^{n+1}\| \le \varepsilon_u\|u^{n+1}\|$ (relative tolerance); stop with an error
        after $i_{\max}$ iterations.
  \item Correctors: $\ddot u^{n+1}$ and $\dot u^{n+1}$ from \eqref{eq:acc_vel}.
  \item Update of the internal variables $\boldsymbol\alpha^{n+1}$ (e.g. $d^{n+1} = \max(d^n, \hat d(u^{n+1}))$ for a damage model).
\end{enumerate}
At $t = 0$ the acceleration is obtained from the equation of motion itself:
$\ddot u^0 = M^{-1}\bigl(f^{\text{ext}}(u^0,0) - C\dot u^0 - f^{\text{int}}(u^0)\bigr)$.

\section{Application: two springs in series, spring 1 damageable}
Node 1 of mass $m_1$ between a damageable spring $k_1^d(u_1) = k_1\bigl(1-d\bigr)$, $d = \max\hat d(u_1)$ and an
elastic spring $k_2$ whose other end is displacement controlled, $u_2(t)$. Then
\begin{equation}
  M = m_1,\quad C = 0,\quad f^{\text{int}}(u_1) = \bigl(k_1^d(u_1)+k_2\bigr)u_1,\quad
  f^{\text{ext}}(t) = k_2\,u_2(t),
\end{equation}
and \eqref{eq:residual}--\eqref{eq:tangent} become
\begin{equation}
  r(u_1) = \frac{m_1}{\beta\Delta t^2}(u_1-\tilde u) + \bigl(k_1^d(u_1)+k_2\bigr)u_1 - k_2u_2^{n+1},
  \qquad
  \frac{\partial r}{\partial u_1} = \frac{m_1}{\beta\Delta t^2} + k_1^d + u_1\frac{\mathrm dk_1^d}{\mathrm du_1} + k_2 .
\end{equation}
The static problem is the limit $m_1\to0$: $R(u_1) = (k_1^d+k_2)u_1 - k_2u_2$, with tangent
$k_t = k_1^d + u_1\,\mathrm dk_1^d/\mathrm du_1 + k_2$, which becomes negative in the softening
range when $k_2 < \max(-\mathrm dF_1/\mathrm du_1)$. The inertia term $m_1/(\beta\Delta t^2)$ added
to the tangent is what restores a strictly increasing residual and a unique root.

\end{document}
